\documentclass[12pt,a4paper,oneside]{article}
\usepackage[brazil]{babel}
\usepackage[utf8]{inputenc}
\usepackage{olymp}
\usepackage{color}
\usepackage[dvipsnames]{xcolor}
\usepackage{listings}

\setcounter{problem}{4}

\begin{document}

\begin{problem}{Conta bits}{bits}{2}

Faça uma função recursiva que conta o número de bits 1 da representação binária de um inteiro $n$. Por exempo, para $n = 14$ a função deve retornar $3$, pois 14 em binário é \texttt{1110}; para $n=2018$ a resposta é 7, pois sua representação em binário é \texttt{11111100010}.

Faça um programa para ler uma lista de números inteiros e, usando a função recursiva descrita acima, informar quantos bits 1 cada um deles possui em binário.

Lembre-se que a representação binária pode ser obtida pelos restos de sucessivas divisões por 2.

\begin{center}
\includegraphics[scale=0.3]{images/exercise_109_0}
\end{center}

%\Notes
%
%Observações, se tiver.

\InputFile

A entrada contém vários casos de teste. Cada caso de teste é uma linha contendo um número inteiro $n$. Restrição: $0 < n \leq 2^{31}-1$.

O valor $0$ encerra a lista de casos de teste e não deve ser processado.

\OutputFile

Seu programa deve gerar uma linha de saída para cada caso de teste, contendo o número de bits 1 da representação binária de $n$.

%\Notes
%Nesta questão, seu programa deve usar a função \texttt{fatorial} dada %acima exatamente como está, não a modifique! Sua
%tarefa é apenas criar o programa com a função \texttt{main}.

\Example

\begin{example}
\exmp{
14
15
2018
1000000000
0
}{
3
4
7
13
}%
\end{example}

%Comentários sobre o exemplo, se necessário.

\end{problem}

\end{document}
